\documentclass{article}
\usepackage{decomp}
\begin{document}
\sourcefile{41467\_2015\_BFncomms8923\_MOESM728\_ESM.pdf}
\decompnote{Supplementary Information of Hua, Laserstein \& Helmstaedter, Nat. Commun. 6:7923 (2015), 5 pages: one supplementary figure on each of pp.\,1–3 (each printed with only the label ``Hua, Laserstein \& Helmstaedter / Suppl. Fig. N'' below it; there are no figure captions in this file), Supplementary Table 1 on p.\,4 and Supplementary Table 2 on p.\,5, each followed by its title and caption. No running headers or page numbers.}

\pdfpage{1}
\begin{center}
\includegraphics[width=0.8\linewidth,height=0.6\textheight,keepaspectratio]{figures/sfig1.png}
\end{center}
\begin{center}
Hua, Laserstein \& Helmstaedter\\
Suppl. Fig. 1
\end{center}
\alttext{Supplementary Fig. 1 is a full page of six square greyscale scanning electron micrographs in three rows of two, panel letters in grey boxes at the upper left (a, b / c, d / e, f), at decreasing magnification from right to left and bottom to top as a zoom series. (f) Lowest magnification: a light grey square block face, rotated a little, on a dark background, with small dark holes and a dark round spot at the left; a white square outline labelled ``e'' near the centre; white scale bar at lower right. (e) The region of the box in f: grey tissue with many lighter round spots (some with dark centres), a few small white spots, a dark ring-shaped hole at lower left, and two small white square outlines labelled ``d'' (left) and ``b'' (centre); white scale bar at lower right. (b) and (d) Neuropil at medium magnification with dark, continuous membranes outlining many profiles, dark mitochondria and myelinated axons; b has a white square outline labelled ``a'' in the centre, d a white square outline labelled ``c'' in the centre; white scale bar at lower right of d (none visible in b). (a) and (c) High magnification: neurite profiles bounded by dark continuous membranes, profiles filled with small round vesicles, mitochondria with cristae; white scale bars at lower right. No scale-bar values are printed in this figure.}

\pdfpage{2}
\begin{center}
\includegraphics[width=0.9\linewidth,height=0.4\textheight,keepaspectratio]{figures/sfig2.png}
\end{center}
\begin{center}
Hua, Laserstein \& Helmstaedter\\
Suppl. Fig. 2
\end{center}
\alttext{Supplementary Fig. 2: two greyscale scanning electron micrographs of neuropil side by side at the top of the page, panel letters ``a'' and ``b'' in grey boxes at upper left, each with a white scale bar at lower right (no values printed). (a) High contrast: dark, sharp membranes around densely packed profiles, many dark mitochondria, several thick black myelin rings (lower left, centre, right) and a large pale blood vessel lumen at the upper left. (b) The same kind of neuropil with clearly lower contrast: the background is mid-grey and the membranes paler grey and less sharp, while the few myelin profiles (an elongated one at upper left, rings at right and lower left) are still dark. The main text (Methods) states that the image in Supplementary Fig. 2b was taken from a block stained without the 120 min 50 °C incubation in UA.}

\pdfpage{3}
\begin{center}
\includegraphics[width=0.95\linewidth,height=0.6\textheight,keepaspectratio]{figures/sfig3.png}
\end{center}
\begin{center}
Hua Laserstein Helmstaedter\\
Suppl. Fig. 3
\end{center}
\alttext{Supplementary Fig. 3: a grid of 14 square greyscale scanning electron micrographs, four rows (a, b, c, d) with labels in dark boxes at the upper left of each image. The first column shows the low-magnification overview of each sample (panels a, b, c, d): the curved sample edge at the upper left (b: a black region with a light grey irregular band at the edge; d: a full round sample face with a dark speckled region covering its lower left part) and homogeneous grey tissue with lighter round spots; each overview carries small white squares labelled a1, a2, a3 (from upper left towards lower right, i.e. from the edge inwards), b1–b3, c1–c3, and d1, d2, and a white scale bar at the bottom (no values printed). The other columns show those regions at high magnification (a1–a3, b1–b3, c1–c3, d1–d2; the fourth column of row d is empty); a white scale bar at the bottom of a1. All high-magnification images show neuropil with dark, continuous membranes, mitochondria and black myelin rings (large thick ones in a2 and a3). c1 is much darker than the rest, with very dark, thick membranes and dark cytoplasm; b1–b3 and d1–d2 have somewhat lower contrast than a1–a3.}

\pdfpage{4}
\begin{center}
\scriptsize
\begin{tabular}{|>{\raggedright\arraybackslash}p{1.5cm}|>{\raggedright\arraybackslash}p{1.8cm}|>{\raggedright\arraybackslash}p{1.8cm}|>{\raggedright\arraybackslash}p{1.8cm}|>{\raggedright\arraybackslash}p{1.6cm}|>{\raggedright\arraybackslash}p{1.8cm}|>{\raggedright\arraybackslash}p{1.8cm}|}
\hline
 & \multicolumn{6}{c|}{\textbf{staining protocol}} \\ \hline
\textbf{incubation steps} & \textbf{Briggman (Fig. 1b)} & \textbf{Holcomb (Fig. 1c)} & \textbf{Mikula (Fig. 1d)} & \textbf{OTO} & \textbf{OPTO (Fig. 3g)} & \textbf{Hua (Fig. 1e)} \\ \hline
\textbf{1.} & 2\% OsO4, 2.5\% ferrocyanide, 0.15 M Cac, pH 7.4 & 2\% OsO4, 2 mM CaCl2 1.5\% ferrocyanide, 0.15 M Cac, pH 7.4 & 1 \% periodic acide, 0.1 M Cac, pH 7.4 & 2\% OsO4, unbuffered & 2\% OsO4, unbuffered & 2\% OsO4, 0.15 M Cac, pH 7.4 \\
 & 1.5 h @ rt & 1.5 h @ rt & 1.5 h @ 4 ºC & 1.5 h @ rt & 1.5 h @ rt & 1.5 h @ rt \\ \hline
 & \multicolumn{6}{c|}{No wash} \\ \hline
\textbf{2} & & & & & 1 \% periodic acide, 0.15 M Cac, pH 7.4 & 2.5\% ferrocyanide, 0.15 M Cac, pH 7.4 \\
 & & & & & 1.5 h @ 4 ºC & 1.5 h @ rt \\ \hline
 & \multicolumn{6}{c|}{0.5 h wash in water x 2} \\ \hline
\textbf{3} & 1\% TCH, unbuffered & 1\% TCH, unbuffered & 1\% TCH, 0.1 M Cac, pH 7.4 & 1\% TCH, unbuffered & 1\% TCH, 0.1 M Cac, pH 7.4 & 1\% TCH, unbuffered \\
 & 0.75 h @ 50 ºC & 0.75 h @ rt & 0.75 h @ 50 ºC & 0.75 h @ rt & 0.75 h @ 50 ºC & 0.75 h @ 40 ºC \\ \hline
 & \multicolumn{6}{c|}{0.5 h wash in water x 2} \\ \hline
\textbf{4} & 2\% OsO4, unbuffered & 2\% OsO4, unbuffered & 2\% OsO4, unbuffered & 2\% OsO4, unbuffered & 2\% OsO4, unbuffered & 2\% OsO4, unbuffered \\
 & 1.5 h @ rt & 1.5 h @ rt & 1.5 h @ rt & 1.5 h @ rt & 1.5 h @ rt & 1.5 h @ rt \\ \hline
 & \multicolumn{6}{c|}{0.5 h wash in water x 2} \\ \hline
\textbf{5} & 1 \% uranyl acetate, unbuffered & 1\% uranyl acetate, unbuffered & 1\% uranyl acetate, unbuffered & 1 \% uranyl acetate, unbuffered & 1 \% uranyl acetate, unbuffered & 1\% uranyl acetate, unbuffered \\
 & 2 h @ 50 ºC & overnight @ 4 ºC & 2 h @ 50 ºC & 2 h @ 50 ºC & 2 h @ 50 ºC & overnight @ 4 ºC, 2 h @ 50 ºC \\ \hline
 & \multicolumn{6}{c|}{0.5 h wash in water x 2} \\ \hline
\textbf{6} & Lead aspartate, pH 5.0 & Lead aspartate, pH 5.5 & Lead aspartate, pH 5.0 & Lead aspartate, pH 5.0 & Lead aspartate, pH 5.0 & Lead aspartate, pH 5.0 \\
 & 2 h @ 50 ºC & 2 h @ 60 ºC & 2 h @ 50 ºC & 2 h @ 50 ºC & 2 h @ 50 ºC & 2 h @ 50 ºC \\ \hline
 & \multicolumn{6}{c|}{0.5 h wash in water x 2} \\ \hline
 & \multicolumn{6}{c|}{dehydration, infiltration and embedding} \\ \hline
\end{tabular}
\end{center}
\decompnote{Table set as printed: the left column gives the step number, each protocol cell holds the solution (upper line) and the time and temperature (lower line); the wash rows span all six protocol columns. ``OsO4'' and ``CaCl2'' are printed without subscripts, temperatures with the ordinal sign ``º'' instead of a degree sign, and ``periodic acide'' as spelled. Also in tables/stable1.csv.}

\noindent Supplementary Table 1

\noindent Detailed comparison of staining steps in employed staining protocols (Fig. 1).

\pdfpage{5}
\begin{center}
\scriptsize
\begin{tabular}{|>{\raggedright\arraybackslash}p{1.5cm}|>{\raggedright\arraybackslash}p{1.6cm}|>{\raggedright\arraybackslash}p{1.3cm}|>{\raggedright\arraybackslash}p{1.3cm}|>{\raggedright\arraybackslash}p{2.6cm}|>{\raggedright\arraybackslash}p{3.0cm}|>{\raggedright\arraybackslash}p{1.5cm}|}
\hline
Experimenter & Species, brain region & Exp. ID & Exp. Date & Variation to standard protocol as in methods & Ultrastructural results & Screening results figure \\ \hline
Hua & Mouse S1 cortex & scYH0009 & 2013/12/11 & rOTO (ref. 18) & Poor penetration, good ultrastructure preservation, good membrane contrast & Fig. 1b \\ \hline
Hua & Mouse S1 cortex & scYH0010 & 2013/12/11 & rOTO (ref. 19) & Poor penetration, good ultrastructure preservation, good membrane contrast & Fig. 1c \\ \hline
Hua & Mouse S1 cortex & scYH0053 & 2014/01/13 & PATCO (ref. 2) & Good penetration, poor ultrastructure preservation, poor membrane contrast, charging & Fig. 1d \\ \hline
Hua & Mouse S1 cortex & scYH0056 & 2014/01/22 & No variation & Penetration, contrast, ultrastructure (all good ) resin hardness (35 nm cut at border, soft core) & Fig. 1e Suppl Fig 2a \\ \hline
Hua \& Laserstein & Mouse S1 cortex & scPL0052 & 2014/08/14 & No accelerator during initial infiltration, Add accelerator on the next day & Penetration, contrast, ultrastructure (all good ) resin hardness (30 nm cut) & Traceability test \\ \hline
Hua & Mouse s1 cortex & scYH1059 & 2014/01/22 & No 2h incubation in UA at 50ºC & Good penetration, good ultrastructure preservation, slightly low membrane contrast & Suppl Fig 2b \\ \hline
Hua & Mouse S1 cortex & scYH1012 & 2015/05/07 & No accelerator during initial infiltration, Added accelerator on the next day instead & Penetration, contrast, ultrastructure (all good ) & Fig. 2 Suppl Fig. 1 \\ \hline
Laserstein & Mouse S1 cortex & scPL0066 & 2014/08/25 & No variation & Penetration, contrast, ultrastructure (all good ) resin hardness (not tested) & Suppl Fig. 3a \\ \hline
Gour & Mouse S1 cortex & scAG0008 & 2014/10/20 & No variation & Penetration, contrast, ultrastructure (all good ) resin hardness (not tested) & Suppl Fig. 3b \\ \hline
Karimi & Mouse PPC & ppcAK0004 & 2014/10/21 & No variation & Penetration, contrast, ultrastructure (all good ), charging effect at the edge area, resin hardness (not tested) & Suppl Fig. 3c \\ \hline
Karimi & Mouse PPC & ppcAK0003 & 2014/10/21 & No variation & Penetration, contrast, ultrastructure (all good ) resin hardness (not tested) & Suppl Fig. 3d \\ \hline
\end{tabular}
\end{center}
\decompnote{Also in tables/stable2.csv. Line breaks inside the cells are not kept. ``Mouse s1 cortex'' (lower-case s) in the scYH1059 row, ``(all good )'' with a space, and ``Suppl Fig'' without a full stop in some cells are as printed.}

\noindent Supplementary Table 2

\noindent Summary of staining experiments reported in the paper for assessing the variability of staining quality over cortex areas and experimentalists.

\end{document}
